\chapter{2011年江西卷（文）}

\section{单选题}

\begin{problem}
若 \((x - \i)\i = y + 2\i\)，\(x\)，\(y \in \R\)，则复数 \(x + y\i =\)（\quad）

\choices
    {\(-2 + \i\)}
    {\(2 + \i\)}
    {\(1 - 2\i\)}
    {\(1 + 2\i\)}
\end{problem}

\begin{answer}
B
\end{answer}

\begin{solution}
展开左边得
\[
(x-\i)\i=x\i-\i^2=1+x\i
\]
比较实部和虚部，得到 \(y=1\)，\(x=2\)，所以 \(x+y\i=2+\i\)，选 B.
\end{solution}

\begin{problem}
若全集 \(U = \{1, 2, 3, 4, 5, 6\}\)，\(M = \{2, 3\}\)，\(N = \{1, 4\}\)，则集合 \(\{5, 6\}\) 等于（\quad）

\choices
    {\(M \cup N\)}
    {\(M \cap N\)}
    {\((\complement_U M) \cup (\complement_U N)\)}
    {\((\complement_U M) \cap (\complement_U N)\)}
\end{problem}

\begin{answer}
D
\end{answer}

\begin{solution}
在全集 \(U\) 中，
\(\complement_U M=\{1,4,5,6\}\)，\(\complement_U N=\{2,3,5,6\}\).
因此
\[
(\complement_U M)\cap(\complement_U N)=\{5,6\}
\]
故选 D.
\end{solution}

\begin{problem}
若 \( f(x) = \frac{1}{\log_{\frac{1}{2}}(2x + 1)} \)，则 \( f(x) \) 的定义域为（\quad）

\choices
    {\(\left(-\frac{1}{2}, 0\right)\)}
    {\(\left(-\frac{1}{2}, +\infty\right)\)}
    {\(\left(-\frac{1}{2}, 0\right) \cup (0, +\infty)\)}
    {\(\left(-\frac{1}{2}, 2\right)\)}
\end{problem}

\begin{answer}
C
\end{answer}

\begin{solution}
对数的真数须为正，且分母不能为零，因此
\(2x+1>0\)，\(\log_{1/2}(2x+1)\ne0\).
这分别给出 \(x>-\frac12\) 和 \(2x+1\ne1\)，即 \(x\ne0\). 所以定义域为
\[
\left(-\frac12,0\right)\cup(0,+\infty)
\]
故选 C.
\end{solution}

\begin{problem}
曲线 \( y = \e^{x} \) 在点 \( A(0,1) \) 处的切线斜率为（\quad）

\choices
    {\(1\)}
    {\(2\)}
    {\(\e\)}
    {\(\frac{1}{\e}\)}
\end{problem}

\begin{answer}
A
\end{answer}

\begin{solution}
函数 \(y=\e^x\) 的导数为 \(y'=\e^x\). 在 \(A(0,1)\) 处，切线斜率为
\[
y'(0)=\e^0=1
\]
故选 A.
\end{solution}

\begin{problem}
设 \(\{a_{n}\}\) 为等差数列，公差 \(d = -2\)，\(S_{n}\) 为其前 \(n\) 项和. 若 \(S_{10} = S_{11}\)，则 \(a_{1} =\)（\quad）

\choices
    {\(18\)}
    {\(20\)}
    {\(22\)}
    {\(24\)}
\end{problem}

\begin{answer}
B
\end{answer}

\begin{solution}
由 \(S_{11}-S_{10}=a_{11}\)，条件 \(S_{10}=S_{11}\) 给出 \(a_{11}=0\). 又
\[
a_{11}=a_1+10d=a_1-20
\]
所以 \(a_1=20\)，故选 B.
\end{solution}

\begin{problem}
观察下列各式： \(7^{2} = 49\)，\(7^{3} = 343\)，\(7^{4} = 2401\)，\(\cdots\)，则 \(7^{2011}\) 的末两位数字为（\quad）

\choices
    {\(01\)}
    {\(43\)}
    {\(07\)}
    {\(49\)}
\end{problem}

\begin{answer}
B
\end{answer}

\begin{solution}
由题给幂的末两位可得 \(7^4\equiv1\pmod {100}\). 因为
\[
2011=4\times502+3
\]
所以
\[
7^{2011}\equiv(7^4)^{502}7^3\equiv343\equiv43\pmod {100}
\]
故末两位为 \(43\)，选 B.
\end{solution}

\begin{problem}
为了普及环保知识，增强环保意识，某大学随机抽取 30 名学生参加环保知识测试，得分 （十分制） 如图所示，假设得分值的中位数为 \(m_{e}\)，众数为 \(m_{o}\)，平均值为 \(\overline{x}\)，则（\quad）

\texfigure[max width=0.72\linewidth]{img_repaint/2011/jiangxi_liberal/q7_fig1.tex}

\choices
    {\(m_{e} = m_{o} = \overline{x}\)}
    {\(m_{e} = m_{o} < \overline{x}\)}
    {\(m_{e} < m_{o} < \overline{x}\)}
    {\(m_{o} < m_{e} < \overline{x}\)}
\end{problem}

\begin{answer}
D
\end{answer}

\begin{solution}
由直方图读得得分 \(3\)，\(4\)，\(5\)，\(6\)，\(7\)，\(8\)，\(9\)，\(10\) 的频数分别为
\(2\)，\(3\)，\(10\)，\(6\)，\(3\)，\(2\)，\(2\)，\(2\).
总人数为 \(30\)，第 \(15\)、\(16\) 个数据分别为 \(5,6\)，故
\[
m_e=\frac{5+6}{2}=\frac{11}{2}
\]
频数最大的得分是 \(5\)，所以 \(m_o=5\). 平均值为
\[
\overline{x}=\frac{3\times2+4\times3+5\times10+6\times6+7\times3+8\times2+9\times2+10\times2}{30}
=\frac{179}{30}>\frac{11}{2}
\]
因此 \(m_o<m_e<\overline{x}\)，选 D.
\end{solution}

\begin{problem}
为了解儿子身高与其父亲身高的关系，随机抽取 5 对父子的身高数据如下：
\begin{center}
\renewcommand{\arraystretch}{1.1}
\begin{tabular}{c|ccccc}
父亲身高 \(x(\mathrm{cm})\) & 174 & 176 & 176 & 176 & 178 \\
\hline
儿子身高 \(y(\mathrm{cm})\) & 175 & 175 & 176 & 177 & 177
\end{tabular}
\end{center}
则 \(y\) 对 \(x\) 的线性回归方程为（\quad）
\choices
    {\(y = x - 1\)}
    {\(y = x + 1\)}
    {\(y = 88 + \frac{1}{2} x\)}
    {\(y = 176\)}
\end{problem}

\begin{answer}
C
\end{answer}

\begin{solution}
由表中数据
\(\overline{x}=\frac{174+176+176+176+178}{5}=176\)，\(\overline{y}=\frac{175+175+176+177+177}{5}=176\).
回归直线斜率为
\[
\hat b=\frac{\sum (x_i-\overline{x})(y_i-\overline{y})}{\sum (x_i-\overline{x})^2}
=\frac{(-2)(-1)+2\cdot1}{(-2)^2+2^2}=\frac12
\]
截距为 \(\overline y-\hat b\,\overline x=176-\frac12\times176=88\)，故回归方程为
\[
y=88+\frac12x
\]
选 C.
\end{solution}

\begin{problem}
将长方体截去一个四棱锥，得到的几何体如图所示，则该几何体的左视图为（\quad）

\bitmapfigure[max width=0.72\linewidth]{img/2011/jiangxi_liberal/q9_fig1.png}

\choices
    {\choicebitmap[width=0.15\paperwidth]{img/2011/jiangxi_liberal/q9_choiceA.png}}
    {\choicebitmap[width=0.15\paperwidth]{img/2011/jiangxi_liberal/q9_choiceB.png}}
    {\choicebitmap[width=0.15\paperwidth]{img/2011/jiangxi_liberal/q9_choiceC.png}}
    {\choicebitmap[width=0.15\paperwidth]{img/2011/jiangxi_liberal/q9_choiceD.png}}
\end{problem}

\begin{answer}
D
\end{answer}

\begin{solution}
沿图示左视方向作正投影，所得外轮廓为正方形. 截去四棱锥后，切割面在左视图中的可见交线连接正方形的左下顶点和右上顶点，形成左下至右上的对角线；其余棱在该方向上的投影不产生新的可见边. 因此左视图为选项 D.
\end{solution}

\begin{problem}
如图，一个“凸轮”放置于直角坐标系 \(x\) 轴上方，其“底端”落在原点 \(O\) 处，一顶点及中心 \(M\) 在 \(y\) 轴正半轴上，它的外围由以正三角形的顶点为圆心，以正三角形的边长为半径的三段等弧组成. 今使“凸轮”沿 \(x\) 轴正向滚动前进，在滚动过程中“凸轮”每时每刻都有一个“最高点”，其中心也在不断移动位置. 则在“凸轮”滚动一周的过程中，将其“最高点”和“中心点”所形成的图形按上，下放置，应大致为（\quad）

\bitmapfigure[max width=0.72\linewidth]{img/2011/jiangxi_liberal/q10_fig1.png}

\choices
    {\choicebitmap[width=0.15\paperwidth]{img/2011/jiangxi_liberal/q10_choiceA.png}}
    {\choicebitmap[width=0.15\paperwidth]{img/2011/jiangxi_liberal/q10_choiceB.png}}
    {\choicebitmap[width=0.15\paperwidth]{img/2011/jiangxi_liberal/q10_choiceC.png}}
    {\choicebitmap[width=0.15\paperwidth]{img/2011/jiangxi_liberal/q10_choiceD.png}}
\end{problem}

\begin{answer}
A
\end{answer}

\begin{solution}
凸轮的边界由三段等弧组成，是等宽曲线；滚动时接触点与最高点的竖直距离始终等于其宽度. 因此最高点的高度保持不变，最高点轨迹为水平直线. 中心 \(M\) 是内接正三角形的中心，不是凸轮的支撑中心，随凸轮转动其高度连续变化，轨迹呈平滑起伏状. 按题意上下放置后，符合这种形状的是选项 A.
\end{solution}

\section{填空题}

\begin{problem}
已知两个单位向量 \(e_{1}\)，\(e_{2}\) 的夹角为 \( \frac{\pi}{3} \)，若向量 \(b_{1} = e_{1} - 2e_{2}\)，\(b_{2} = 3e_{1} + 4e_{2}\)，则 \( b_{1} \cdot b_{2} = \)\fillinblank{}.
\end{problem}

\begin{answer}
\(-6\)
\end{answer}

\begin{solution}
由单位向量夹角为 \(\frac{\pi}{3}\)，有 \(e_1\cdot e_1=e_2\cdot e_2=1\)，且 \(e_1\cdot e_2=\cos\frac{\pi}{3}=\frac12\). 因此
\[
b_1\cdot b_2=(e_1-2e_2)\cdot(3e_1+4e_2)=3+4\cdot\frac12-6\cdot\frac12-8=-6
\]
\end{solution}

\begin{problem}
若双曲线 \(\frac{y^2}{16} - \frac{x^2}{m} = 1\) 的离心率 \(e = 2\)，则 \(m = \)\fillinblank{}.
\end{problem}

\begin{answer}
\(48\)
\end{answer}

\begin{solution}
该双曲线的实半轴为 \(a_h=4\)，虚半轴平方为 \(b_h^2=m\)，故焦半距满足
\[
c_h^2=a_h^2+b_h^2=16+m
\]
由离心率 \(e=\frac{c_h}{a_h}=2\)，得 \(c_h=8\). 因而
\[
16+m=8^2=64
\]
所以 \(m=48\).
\end{solution}

\begin{problem}
下图是某算法程序框图，则程序运行后输出的结果是\fillinblank{}.

\texfigure[max width=0.96\linewidth]{img_repaint/2011/jiangxi_liberal/q13_fig1.tex}
\end{problem}

\begin{answer}
\(6\)
\end{answer}

\begin{solution}
按流程图从 \(s=0\)，\(n=1\) 开始执行. 第一次循环后
\(s=(0+1)\times1=1\)，\(n=2\)
此时 \(s>3\) 不成立，继续循环. 第二次循环后
\(s=(1+2)\times2=6\)，\(n=3\)
此时 \(s>3\) 成立，输出 \(s=6\).
\end{solution}

\begin{problem}
已知角 \(\theta\) 的顶点为坐标原点，始边为 \(x\) 轴的正半轴. 若 \(P(4, y)\) 是角 \(\theta\) 终边上一点，且 \(\sin \theta = -\frac{2\sqrt{5}}{5}\)，则 \(y =\fillinblank{}\).
\end{problem}

\begin{answer}
\(-8\)
\end{answer}

\begin{solution}
点 \(P(4,y)\) 到原点的距离为 \(r=\sqrt{16+y^2}\)，所以
\[
\sin\theta=\frac{y}{r}=-\frac{2\sqrt5}{5}
\]
平方并化简得
\(\frac{y^2}{16+y^2}=\frac45\)，\(\Longrightarrow\)，\(y^2=64\).
由于 \(P\) 的横坐标为正而正弦为负，终边在第四象限，故 \(y<0\)，从而 \(y=-8\).
\end{solution}

\begin{problem}
对于 \(x \in \R\)，不等式 \(|x + 10| - |x - 2| \geqslant 8\) 的解集为\fillinblank{}.
\end{problem}

\begin{answer}
\([0,+\infty)\)
\end{answer}

\begin{solution}
按 \(x=-10\) 和 \(x=2\) 分段讨论：
\[
|x+10|-|x-2|=
\begin{cases}
-12,&x<-10,\\
2x+8,&-10\le x<2,\\
12,&x\ge2.
\end{cases}
\]
第一段不满足不等式；第二段要求 \(2x+8\ge8\)，即 \(x\ge0\)；第三段恒满足. 因此解集为 \([0,+\infty)\).
\end{solution}

\section{解答题}

\begin{problem}
某饮料公司对一名员工进行测试以便确定其考评级别. 公司准备了两种不同的饮料共5杯，其颜色完全相同，并且其中3杯为A饮料，另外2杯为B饮料，公司要求此员工一一品尝后，从5杯饮料中选出3杯A饮料. 若该员工3杯都选对，则评为优秀；若3杯选对2杯，则评为良好；否则评为及格. 假设此人对 \(A\) 和 \(B\) 两种饮料没有鉴别能力.
\begin{enumerate}
\item 求此人被评为优秀的概率；
\item 求此人被评为良好及以上的概率.
\end{enumerate}
\end{problem}

\begin{solution}
\begin{enumerate}
\item 因为没有鉴别能力，任意从 \(5\) 杯中选出 \(3\) 杯的结果等可能，共有 \(\mathrm C_5^3=10\) 种. 三杯都选中 A 饮料只有 \(\mathrm C_3^3\mathrm C_2^0=1\) 种，所以被评为优秀的概率为
\[
P_1=\frac{\mathrm C_3^3\mathrm C_2^0}{\mathrm C_5^3}=\frac1{10}
\]
\item 恰好选中 \(2\) 杯 A 饮料、\(1\) 杯 B 饮料有 \(\mathrm C_3^2\mathrm C_2^1=6\) 种. 良好及以上包括三杯全对和恰好两杯对，因此
\[
P_2=\frac{\mathrm C_3^3\mathrm C_2^0+\mathrm C_3^2\mathrm C_2^1}{\mathrm C_5^3}
=\frac{1+6}{10}=\frac7{10}
\]
\end{enumerate}
\end{solution}

\begin{problem}
在 \(\triangle ABC\) 中，\(A\)，\(B\)，\(C\) 的对边分别是 \(a\)，\(b\)，\(c\). 已知 \(3a \cos A = c \cos B + b \cos C\).
\begin{enumerate}
\item 求 \( \cos A \) 的值；
\item 若 \(a = 1\)，\(\cos B + \cos C = \frac{2\sqrt{3}}{3}\)，求边 \(c\) 的值.
\end{enumerate}
\end{problem}

\begin{solution}
\begin{enumerate}
\item 由余弦定理，
\(c\cos B=\frac{a^2+c^2-b^2}{2a}\)，\(b\cos C=\frac{a^2+b^2-c^2}{2a}\).
两式相加得 \(c\cos B+b\cos C=a\). 代入已知条件，得到 \(3a\cos A=a\). 因为 \(a>0\)，所以
\[
\cos A=\frac13
\]
\item 由上问，\(\cos A=\frac13\)，从而
\[
\sin\frac A2=\sqrt{\frac{1-\cos A}{2}}=\frac1{\sqrt3}
\]
又
\[
\cos B+\cos C
=2\cos\frac{B+C}{2}\cos\frac{B-C}{2}
=2\sin\frac A2\cos\frac{B-C}{2}
\]
结合题设 \(\cos B+\cos C=\frac2{\sqrt3}\)，可得
\[
\cos\frac{B-C}{2}=1
\]
由于 \(B\)，\(C\in(0,\pi)\)，有 \(\frac{B-C}{2}\in(-\frac\pi2,\frac\pi2)\)，故 \(B=C\)，从而 \(b=c\). 再由余弦定理和 \(a=1\)，
\[
1=b^2+c^2-2bc\cos A
=2c^2\left(1-\frac13\right)=\frac43c^2
\]
因此 \(c=\frac{\sqrt3}{2}\).
\end{enumerate}
\end{solution}

\begin{problem}
如图，在 \(\triangle ABC\) 中，\(\angle B = \frac{\pi}{2}\)，\(AB = BC = 2\). \(P\) 为 \(AB\) 边上动一点，\(PD \parallel BC\) 交 \(AC\) 于点 \(D\)，现将 \(\triangle PDA\) 沿 \(PD\) 翻折至 \(\triangle PDA'\)，使平面 \(PDA' \perp\) 平面 \(PBCD\).
\begin{enumerate}
\item 当棱锥 \(A' - PBCD\) 的体积最大时，求 \(PA\) 的长；
\item 若点 \(P\) 为 \(AB\) 的中点，\(E\) 为 \(A'C\) 的中点，求证： \(A'B \perp DE\).
\end{enumerate}

\bitmapfigure[max width=0.52\linewidth]{img/2011/jiangxi_liberal/q18_fig1.png}
\end{problem}

\begin{solution}
\begin{enumerate}
\item 设 \(PA=x\)，则 \(0\le x\le2\). 取
\(B=(0,0,0)\)，\(A=(2,0,0)\)，\(C=(0,2,0)\).
由 \(PD\parallel BC\)，得 \(P=(2-x,0,0)\)，\(D=(2-x,x,0)\). 折叠后可取 \(A'=(2-x,0,x)\)，所以棱锥的高为 \(x\). 又
\(S_{\triangle ABC}=2\)，\(S_{\triangle APD}=\frac{x^2}{2}\)
故底面 \(PBCD\) 的面积为 \(2-\frac{x^2}{2}\). 体积为
\(V(x)=\frac13\left(2-\frac{x^2}{2}\right)x =\frac{x(4-x^2)}6\)，\(0\le x\le2\).
求导得
\[
V'(x)=\frac{4-3x^2}{6}
\]
于是 \(V\) 在 \(x=\frac2{\sqrt3}\) 处由增转减，且 \(V(0)=V(2)=0\)，故体积最大时
\[
PA=x=\frac{2\sqrt3}{3}
\]
\item 当 \(P\) 为 \(AB\) 的中点时，取 \(x=1\)，则
\(A'=(1,0,1)\)，\(B=(0,0,0)\)，\(D=(1,1,0)\)，\(C=(0,2,0)\).
中点 \(E\) 为
\[
E=\frac{A'+C}{2}=\left(\frac12,1,\frac12\right)
\]
因此
\(\overrightarrow{A'B}=(-1,0,-1)\)，\(\overrightarrow{DE}=\left(-\frac12,0,\frac12\right)\)
从而
\[
\overrightarrow{A'B}\cdot\overrightarrow{DE}
=(-1)\left(-\frac12\right)+(-1)\left(\frac12\right)=0
\]
所以 \(A'B\perp DE\).
\end{enumerate}
\end{solution}

\begin{problem}
已知过抛物线 \( y^{2} = 2px \) （\( p > 0 \)） 的焦点，斜率为 \( 2\sqrt{2} \) 的直线交抛物线于 \( A(x_{1}, y_{1}) \)，\( B(x_{2}, y_{2}) \) （\( x_{1} < x_{2} \)） 两点，且 \( |AB| = 9 \).
\begin{enumerate}
\item 求该抛物线的方程；
\item O 为坐标原点，C 为抛物线上一点. 若 \( \overrightarrow{OC} = \overrightarrow{OA} + \lambda \overrightarrow{OB} \)，求 \( \lambda \) 的值.
\end{enumerate}
\end{problem}

\begin{solution}
\begin{enumerate}
\item 抛物线 \(y^2=2px\) 的焦点为 \((\frac p2,0)\)，故所给直线为
\[
y=2\sqrt2\left(x-\frac p2\right)
\]
代入抛物线方程，得
\(8\left(x-\frac p2\right)^2=2px\)，\(\Longrightarrow\)，\(4x^2-5px+p^2=0\).
所以
\(x_1=\frac p4\)，\(x_2=p\).
直线的斜率为 \(2\sqrt2\)，故
\[
|AB|=\sqrt{1+(2\sqrt2)^2}\,(x_2-x_1)
=3\cdot\frac{3p}{4}=\frac{9p}{4}
\]
由 \(|AB|=9\) 得 \(p=4\)，抛物线方程为
\[
y^2=8x
\]
\item 此时焦点为 \((2,0)\)，直线与抛物线的交点为
\(A=(1,-2\sqrt2)\)，\(B=(4,4\sqrt2)\).
由向量条件，点 \(C\) 的坐标为
\[
C=A+\lambda B=(1+4\lambda,-2\sqrt2+4\sqrt2\lambda)
\]
代入 \(y^2=8x\)，得
\(8(2\lambda-1)^2=8(1+4\lambda)\)，\(\Longrightarrow\)，\(\lambda(\lambda-2)=0\).
故 \(\lambda=0\) 或 \(\lambda=2\). 其中 \(\lambda=0\) 对应 \(C=A\)，\(\lambda=2\) 对应 \(C=(9,6\sqrt2)\)；题面未排除 \(C=A\)，所以两值都应保留.
\end{enumerate}
\end{solution}

\begin{problem}
设 \( f(x) = \frac{1}{3} x^3 + mx^2 + nx \).
\begin{enumerate}
\item 如果 \( g(x) = f'(x) - 2x - 3 \) 在 \( x = -2 \) 处取得最小值 \(-5\)，求 \( f(x) \) 的解析式；
\item 如果 \(m + n < 10\) （\(m\)，\(n \in \N_+\)），\(f(x)\) 的单调递减区间的长度是正整数，试求 \(m\) 和 \(n\) 的值. （注：区间 \((a, b)\) 的长度为 \(b - a\)）
\end{enumerate}
\end{problem}

\begin{solution}
\begin{enumerate}
\item
\[
f'(x)=x^2+2mx+n
\]
所以
\[
g(x)=x^2+(2m-2)x+n-3
\]
这是开口向上的二次函数，其对称轴为 \(x=1-m\). 由最小值在 \(x=-2\) 处取得，得 \(1-m=-2\)，即 \(m=3\). 再由
\[
g(-2)=4-8+n-3=n-7=-5
\]
得 \(n=2\). 因此
\[
f(x)=\frac13x^3+3x^2+2x
\]
\item 此时
\[
f'(x)=x^2+2mx+n
\]
若 \(m^2\le n\)，则 \(f'(x)=(x+m)^2+n-m^2\ge0\)，不存在长度为正的单调递减区间. 因而必须有 \(m^2>n\). 此时 \(f'(x)<0\) 的区间长度为
\[
2\sqrt{m^2-n}
\]
设该长度为 \(L\in\N_+\)，则 \(L^2=4(m^2-n)\). 因 \(m^2-n\) 为整数，\(L^2\) 是 \(4\) 的倍数，故 \(L\) 为偶数. 令 \(k=L/2\)，便有 \(m^2-n=k^2\)，其中 \(k\) 为正整数，即
\(n=m^2-k^2\)，\(1\le k\le m-1\).
由于 \(n\ge1\) 且 \(m^2>n\)，\(m=1\) 不可能，故 \(m\ge2\). 结合 \(m+n<10\) 枚举：\(m=2\)，\(k=1\) 时 \(n=3\)；\(m=3\)，\(k=1\) 时 \(n=8\) 但 \(m+n=11\) 不合题意；\(m=3\)，\(k=2\) 时 \(n=5\). 当 \(m\ge4\) 时，
\[
m+n\ge m+(2m-1)=3m-1\ge11
\]
也不合题意. 因此
\[
(m,n)=(2,3)\quad\text{或}\quad(m,n)=(3,5)
\]
\end{enumerate}
\end{solution}

\begin{problem}
\begin{enumerate}
\item 已知两个等比数列 \(\{a_{n}\}\)，\(\{b_{n}\}\)，满足 \(a_{1} = a\)（\(a > 0\)），\(b_{1} - a_{1} = 1\)，\(b_{2} - a_{2} = 2\)，\(b_{3} - a_{3} = 3\)，若数列 \(\{a_{n}\}\) 唯一，求 \(a\) 的值；
\item 是否存在两个等比数列 \(\{a_{n}\}\)，\(\{b_{n}\}\)，使得 \(b_{1} - a_{1}\)，\(b_{2} - a_{2}\)，\(b_{3} - a_{3}\)，\(b_{4} - a_{4}\) 成公差不为 0 的等差数列？ 若存在，求 \(\{a_{n}\}\)，\(\{b_{n}\}\) 的通项公式；若不存在，说明理由.
\end{enumerate}
\end{problem}

\begin{solution}
\begin{enumerate}
\item 设两个等比数列的首项和公比分别为 \(a\)，\(q\) 及 \(b\)，\(r\). 按等比数列的定义，公比不为 \(0\). 由 \(b-a=1\) 得 \(b=a+1\)，再由第二、三项的差分别为 \(2,3\)，得
\((a+1)r-aq=2\)，\((a+1)r^2-aq^2=3\).
由第一式 \(r=\frac{2+aq}{a+1}\)，代入第二式并整理，得到
\[
a(q^2-4q+3)=1
\]
即 \(q\) 是方程
\[
a q^2-4a q+3a-1=0
\]
的根. 其判别式为 \(\Delta=16a^2-4a(3a-1)=4a(a+1)>0\)，所以对每个 \(a>0\) 都有两个不同的实根. 对任一根，第一式唯一确定 \(r\)，且 \(r\ne0\)：若 \(r=0\)，则由第一式得 \(q=-\frac2a\)，代入第二式将得到 \(-\frac4a=3\)，矛盾.

当 \(a\ne\frac13\) 时，方程的两个根均不为 \(0\)，因而均对应合法的公比，得到两个不同的数列 \(\{a_n\}\)，不满足唯一性. 当 \(a=\frac13\) 时，方程化为 \(q(q-4)=0\)，其中 \(q=0\) 不合等比数列公比的定义，故唯一可取 \(q=4\). 此时 \(b=\frac43\)，由第一式得 \(r=\frac52\)，所以
\[
a=\frac13
\]
\item 记
\[
d_n=b_n-a_n=br^{n-1}-aq^{n-1}
\]
若 \(d_1\)，\(d_2\)，\(d_3\)，\(d_4\) 成等差数列，则两个二阶差分均为零，即
\[
b(r-1)^2-a(q-1)^2=0
\]
\[
br(r-1)^2-aq(q-1)^2=0
\]
令 \(X=b(r-1)^2\)，\(Y=a(q-1)^2\)，则 \(X=Y\)，且 \(rX=qY\)，从而
\[
(r-q)X=0
\]
若 \(X\ne0\)，则 \(r=q\)，于是 \(d_n=(b-a)q^{n-1}\). 由于等差数列的公差非零，\(b-a\ne0\)，由其前三项二阶差分为零得 \((q-1)^2=0\)，即 \(q=1\)，这又使 \(d_n\) 为常数，公差为零，矛盾.
若 \(X=0\)，则 \(Y=0\). 当 \(a\)，\(b\ne0\) 时，\(q=r=1\)，同样使 \(d_n\) 为常数. 若 \(a=0\)，则 \(d_n=br^{n-1}\)；当 \(b\ne0\) 时，等差条件的二阶差分给出 \(b(r-1)^2=0\)，故 \(r=1\)，公差仍为零；当 \(b=0\) 时四项全为零. \(b=0\) 的情形同理. 因此不存在满足条件的两个等比数列.
\end{enumerate}
\end{solution}
